% ---------------------------------------------------------------------------
% Template for a single-claim submission.
%
% revtex4-2 is the default because most of the target categories (hep-ph,
% astro-ph.CO, gr-qc) read as APS-style. Swap to `article` for math-ph.
%
% BEFORE WRITING ANY PROSE, fill in the abstract. If the abstract cannot be
% written without pointing at another document, the paper's boundary is wrong
% and the scope needs cutting, not explaining.
% ---------------------------------------------------------------------------
\documentclass[aps,prd,twocolumn,superscriptaddress,nofootinbib,showpacs]{revtex4-2}

\usepackage{amsmath,amssymb}
\usepackage{graphicx}
\usepackage[colorlinks=true,citecolor=blue,linkcolor=blue,urlcolor=blue]{hyperref}
\usepackage{bm}

\begin{document}

\title{TITLE --- state the one claim, not the programme}

\author{Justin Pulford}
\email{pulfordj420@gmail.com}
\affiliation{Unaffiliated}

\date{\today}

% ---------------------------------------------------------------------------
% ABSTRACT. Hard limit 1920 characters in arXiv metadata; aim well under.
% A moderator reads this and nothing else. It must contain, in order:
%   (1) the setting, in one sentence a non-specialist in the subfield follows;
%   (2) the specific claim, with its number;
%   (3) how it is obtained;
%   (4) what would falsify it;
%   (5) what is assumed and not derived --- stated here, not buried.
% Point (5) is not a weakness to hide. A paper that names its own assumptions
% in the abstract reads as competent; one that does not reads as overclaiming,
% and overclaiming is what draws reclassification.
% ---------------------------------------------------------------------------
\begin{abstract}
ABSTRACT.
\end{abstract}

\maketitle

% ---------------------------------------------------------------------------
\section{Introduction}
% What is already known, with citations. Where this differs. Why the difference
% is testable. Do not open with the framework --- open with the problem the
% field already recognises.
% ---------------------------------------------------------------------------

\section{Setup}
% Definitions and notation. Standard symbols only. Anything imported from other
% work is stated with its own citation and its own uncertainty, never assumed
% available to the reader.

\section{Result}
% The derivation. Every step reproducible from what is on the page.

\section{What is assumed}
% A short explicit section. Which inputs are derived, which are assumed, which
% are selected by data, and which are fitted. This is the corpus's strongest
% habit and it transfers directly --- but phrase it as ordinary scientific
% hedging, not as an internal grading vocabulary.

\section{Falsification}
% The number or observation that kills the claim, and the precision needed.

\section{Conclusion}

\begin{acknowledgments}
\end{acknowledgments}

\bibliographystyle{apsrev4-2}
\bibliography{refs}

\end{document}
